\subsection{A net hedge still needs its gross transactions}
\label{sec:net-hedge-timing}

A company expecting dollar sales and dollar expenses may manage only their
net currency exposure. Hedge accounting still needs both gross amounts and the
periods in which they affect earnings. IFRS~9.6.6.1 and B6.6.1--8 require
individually eligible items and actual group risk management. For offsetting
cash-flow positions, the designated risk must be foreign currency risk, and
the designation must specify nature, volume and the relevant reporting periods.
An abstract label such as ``net dollars 20'' is insufficient
\citep[IFRS~9.6.6.1--3, B6.6.1--8]{eu20231803}.

Use a euro-functional company with highly probable external cash sales of
USD~100 and payments of USD~80 for services consumed immediately. The
designation identifies those transactions, their earnings periods and a forward
to sell USD~20 for EUR~20. Assume actual net risk management, an eligible
external derivative, contemporaneous documentation and all effectiveness
conditions. Interest, discount, forward points and credit effects are zero in
this illustration; the whole forward is designated. At inception the rate is
EUR~1 per dollar and the derivative is worth zero. At execution the rate is
EUR~1.10, so sales are EUR~110, service expense EUR~88 and the forward loses
EUR~2. No inventory asset is created in this example.

The effectiveness comparison retains the gross changes: the expense's
currency effect of minus 8 and the forward loss of 2 offset the sales effect
of plus 10. IFRS~9.B6.6.5 and B6.6.9 require the offsetting item to enter
that comparison alongside the instrument, without recognizing a forecast
transaction before it occurs. When both transactions affect the same period,
recognize the effective forward loss of 2 in OCI and subsequently release it
as a separate hedge expense of 2. Revenue remains 110 and service expense 88;
profit is $110-88-2=20$. The hedge result is presented separately because the
offsetting risks affect different statement lines
\citep[IFRS~9.6.5.11, 6.6.4, B6.6.5, B6.6.9, B6.6.15]{eu20231803}.
Showing revenue 108 would preserve profit but misstate the required
presentation. Opening cash and equity 200 become 220 after gross operating
receipts 110, service payments 88 and forward settlement 2. The reserve and
derivative liability both close at zero. The designated forward follows the
operating classification of the hedged flows \citep[IAS~7.16]{eu20231803}.

Move only the service expense into the next period, retaining the same
EUR~1.10 rate and a forward that settles with that expense. The designation
already identifies these different earnings periods. At the first close the
company has cash revenue 110 and a forward liability of 2. The revenue line
must still use its transaction-date spot rate. To reflect the net hedge's
timing, recognize a separate hedge expense of 10 when the sale occurs, with
the corresponding credit to the cash-flow hedge reserve. Together with the
forward's OCI loss of 2, the reserve is a credit of 8. First-period profit is
100 and OCI is 8; cash 310 less the derivative liability 2 equals equity 308.
The 8 reserve carries the offset relating to the later expense; it is neither
an additional derivative asset nor recognized service income
\citep[IFRS~9.B6.6.9, B6.6.15]{eu20231803}.

In period 2, record service expense and payment 88, settle the forward
liability for 2, and release the reserve 8 as separate hedge income. Profit
is minus 80 and OCI minus 8. Cash falls by 90 to 220, which also equals
$308-80-8$. Across both periods profit remains 20 and OCI nets to zero.
Releasing only the derivative loss when the sale occurs would give first-period
profit 108 and miss this timing rule. A later rate change would require new
measurements of all relevant components; the constant second-period rate
isolates the recognition issue without supplying a general valuation model.

With sales and service payments both USD~100, the position is nil and no
derivative is needed. Such a designation is allowed only under all four
conditions in IFRS~9.6.6.6: an actual rolling net-risk strategy, changes in net
size over that strategy's life, normal use of eligible instruments and hedge
accounting when the position is nonnil, and otherwise inconsistent recognition
of the offsetting risks. Conditional on those facts, the earlier-sale example
has revenue 110, separate hedge expense 10 and reserve 10 in period 1.
Expense 110 and separate hedge income 10 in period 2 leave lifetime profit
and cash flow zero. Equal gross totals alone do not establish eligibility
\citep[IFRS~9.6.6.6, B6.6.3, B6.6.6, B6.6.10]{eu20231803}.

The group record therefore preserves each gross exposure's identity, reporting
entity, currency, risk, amount, recognition period and admitted valuation, plus
the strategy, designation, instrument, effectiveness calculation and reserve
history. A layer also needs an identified and tracked overall pool and the
conditions in 6.6.3. A group without offsetting risks apportions the instrument's
net gain or loss systematically to the affected lines without grossing it up;
a fair-value hedge adjusts each relevant asset or liability separately
\citep[IFRS~9.6.6.3--5, B6.6.11--16]{eu20231803}. A group of items is not
permission to ignore consolidated-group external-party requirements: retain
6.2.3 and 6.3.5--6's distinct tests. These methods need explicit successor
records and complete exchanges; the existing closed financing profiles are
preserved until those additions are specified.
